\documentclass{article}
\usepackage{amsmath,amssymb,amsthm}
\usepackage{algorithm}
\usepackage{algpseudocode}
\usepackage{enumitem}
\usepackage{hyperref}
\newtheorem{theorem}{Theorem}
\newtheorem{lemma}{Lemma}
\newtheorem{assumption}{Assumption}
\newcommand{\R}{\mathbb{R}}
\newcommand{\norm}[1]{\left\lVert#1\right\rVert}
\newcommand{\inner}[2]{\left\langle #1,#2\right\rangle}
\begin{document}

\section*{Algorithm Name}
\textbf{Heavy–Ball Method (HB)}  

The method augments the gradient descent step with a momentum term that uses the previous displacement.

\section*{Mathematical Formulation}
Let $f:\R^{d}\rightarrow\R$ be differentiable.  
Given step size $\alpha>0$ and momentum parameter $\beta\in[0,1)$, the iterates $\{x_k\}_{k\ge 0}$ are defined by  

\[
\boxed{
\begin{aligned}
x_{k+1}&=x_k+\beta\bigl(x_k-x_{k-1}\bigr)-\alpha\,\nabla f(x_k),\qquad k\ge 1,\\[2mm]
x_1&=x_0-\alpha\,\nabla f(x_0) \quad\text{(first step, no momentum).}
\end{aligned}}
\]

Denote the optimal point $x^{\star}=\arg\min f(x)$ (unique under strong convexity).

\section*{Pseudocode}
\begin{algorithm}[H]
\caption{Heavy–Ball Method}
\begin{algorithmic}[1]
\Require Initial point $x_0\in\R^{d}$, step size $\alpha>0$, momentum $\beta\in[0,1)$, max iterations $T$.
\State $x_{-1}\gets x_0$ \Comment{dummy previous iterate}
\State $g_0\gets \nabla f(x_0)$
\State $x_1\gets x_0-\alpha g_0$
\For{$k=1$ \textbf{to} $T-1$}
    \State $g_k\gets \nabla f(x_k)$
    \State $x_{k+1}\gets x_k+\beta\,(x_k-x_{k-1})-\alpha\,g_k$
\EndFor
\State \Return $x_T$
\end{algorithmic}
\end{algorithm}

Termination may also be based on $\norm{\nabla f(x_k)}\le\varepsilon$ or on a relative decrease of $f$.

\section*{Dry Run Example}
Consider the one‑dimensional quadratic  

\[
f(w)=(w-3)^2,\qquad \nabla f(w)=2(w-3).
\]

Choose $\alpha=0.1$, $\beta=0.5$, $w_0=0$.

\begin{center}
\begin{tabular}{c|c|c|c}
$k$ & $w_k$ & $\nabla f(w_k)$ & $f(w_k)$\\\hline
0 & $0$ & $2(0-3)=-6$ & $9$\\
1 & $w_1=w_0-\alpha\nabla f(w_0)=0-0.1(-6)=0.6$ & $2(0.6-3)=-4.8$ & $(0.6-3)^2=5.76$\\
2 & $w_2=w_1+\beta(w_1-w_0)-\alpha\nabla f(w_1)=0.6+0.5(0.6-0)-0.1(-4.8)=0.6+0.3+0.48=1.38$ & $2(1.38-3)=-3.24$ & $(1.38-3)^2=2.6244$\\
3 & $w_3=w_2+\beta(w_2-w_1)-\alpha\nabla f(w_2)=1.38+0.5(1.38-0.6)-0.1(-3.24)=1.38+0.39+0.324=2.094$ & $2(2.094-3)=-1.812$ & $(2.094-3)^2=0.821\,.$
\end{tabular}
\end{center}

The iterates quickly approach the minimiser $w^{\star}=3$.

\newpage
\section*{Assumptions}
\begin{assumption}[Strong convexity and smoothness]\label{ass:sc-smooth}
The function $f$ satisfies, for some constants $0<\mu\le L$,
\begin{align}
\frac{\mu}{2}\norm{x-y}^{2}&\le f(y)-f(x)-\inner{\nabla f(x)}{y-x}\le\frac{L}{2}\norm{x-y}^{2},
\quad\forall x,y\in\R^{d}. \tag{A1}
\end{align}
Consequences:
\begin{enumerate}[label=(\roman*)]
\item $\mu$‑strong convexity: $f(y)\ge f(x)+\inner{\nabla f(x)}{y-x}+\frac{\mu}{2}\norm{y-x}^{2}$.
\item $L$‑smoothness: $\norm{\nabla f(x)-\nabla f(y)}\le L\norm{x-y}$.
\end{enumerate}
\end{assumption}
\begin{assumption}[Parameter region]\label{ass:params}
The stepsize $\alpha$ and momentum $\beta$ satisfy
\begin{align}
0<\alpha<\frac{2}{L+\mu},\qquad
0\le\beta<\left(\frac{1-\sqrt{\alpha\mu}}{1+\sqrt{\alpha\mu}}\right)^{2}. \tag{A2}
\end{align}
Under \eqref{ass:params} the spectral radius of the iteration matrix (defined below) is strictly smaller than one.
\end{assumption}

\section*{Main Convergence Theorem}
\begin{theorem}[Linear convergence of Heavy–Ball]\label{thm:hb}
Assume \ref{ass:sc-smooth} and \ref{ass:params}.  
Let $\{x_k\}$ be generated by the Heavy–Ball method. Define the Lyapunov function  

\[
\Phi_k:=f(x_k)-f(x^{\star})+\frac{\gamma}{2}\norm{x_k-x_{k-1}}^{2},
\qquad\gamma:=\frac{1-\beta}{\alpha}-\frac{L}{2}>0 .
\]

Then there exists a constant $\rho\in(0,1)$, depending only on $\alpha,\beta,\mu,L$, such that  

\[
\Phi_{k+1}\le\rho\,\Phi_k,\qquad\forall k\ge 1.
\]

Consequently,
\[
f(x_k)-f(x^{\star})\le \Phi_k\le \rho^{k-1}\Phi_1,
\qquad 
\norm{x_k-x^{\star}}\le C\,\rho^{k/2},
\]
for a constant $C>0$ independent of $k$.
\end{theorem}

\section*{Theoretical Properties}
\begin{itemize}
\item \textbf{Stability.} Condition \eqref{ass:params} guarantees that the iteration matrix has eigenvalues inside the unit disc; thus the method is BIBO‑stable for quadratic functions.
\item \textbf{Hyperparameter sensitivity.} The admissible region shrinks as the condition number $\kappa=L/\mu$ grows. For ill‑conditioned problems one typically chooses $\alpha\approx 4/(L+\mu)$ and $\beta\approx\bigl(\frac{\sqrt{\kappa}-1}{\sqrt{\kappa}+1}\bigr)^{2}$, the optimal values that minimise $\rho$.
\item \textbf{Comparison to baselines.} Gradient descent with stepsize $2/(L+\mu)$ yields rate $(\frac{L-\mu}{L+\mu})^{k}$, whereas Heavy–Ball can achieve the accelerated rate $(\frac{\sqrt{\kappa}-1}{\sqrt{\kappa}+1})^{k}$, strictly faster when $\kappa>1$.
\end{itemize}

\section*{Discussion}
The theorem shows that, under the classical strong‑convexity/smoothness regime, the Heavy–Ball method enjoys a deterministic linear convergence guarantee. The proof relies on a carefully crafted Lyapunov function that couples function suboptimality with the kinetic energy $\norm{x_k-x_{k-1}}^{2}$. The parameter region \eqref{ass:params} is both necessary and sufficient for the quadratic case; for general smooth strongly convex functions the same region suffices because of the sandwich inequality \eqref{ass:sc-smooth}.

\newpage
\section*{Preliminaries (Definitions and Lemmas)}
\begin{lemma}[Three‑point identity]\label{lem:three}
For any $x,y\in\R^{d}$,
\[
\inner{\nabla f(x)-\nabla f(y)}{x-y}
= \norm{\nabla f(x)-\nabla f(y)}^{2}
+\inner{\nabla f(x)-\nabla f(y)}{x-y-\bigl(\nabla f(x)-\nabla f(y)\bigr)} .
\]
\end{lemma}
\begin{proof}
Straightforward expansion of the inner product.
\end{proof}

\begin{lemma}[Descent lemma]\label{lem:descent}
If $f$ is $L$‑smooth, then for all $x,y$,
\[
f(y)\le f(x)+\inner{\nabla f(x)}{y-x}+\frac{L}{2}\norm{y-x}^{2}.
\]
\end{lemma}
\begin{proof}
Standard consequence of $L$‑smoothness; see e.g. Nesterov (2004).
\end{proof}

\section*{Main Proof (Step‑by‑Step Derivation)}
Let $e_k:=x_k-x^{\star}$ and $s_k:=x_k-x_{k-1}$.  
Recall that $\nabla f(x^{\star})=0$.

\begin{enumerate}[label={[\textbf{Step \arabic*}]}]
\item \textbf{Error recursion.}  
Using the update rule,
\[
\begin{aligned}
e_{k+1}
&=x_{k+1}-x^{\star}\\
&=x_k+\beta s_k-\alpha\nabla f(x_k)-x^{\star}\\
&=e_k+\beta s_k-\alpha\bigl(\nabla f(x_k)-\nabla f(x^{\star})\bigr). \tag{1}
\end{aligned}
\]

\item \textbf{Relation between $s_{k+1}$ and $e_{k+1},e_k$.}  
By definition $s_{k+1}=x_{k+1}-x_k=e_{k+1}-e_k$. Hence
\[
s_{k+1}=e_{k+1}-e_k. \tag{2}
\]

\item \textbf{Lyapunov candidate.}  
Define
\[
\Phi_k:=f(x_k)-f(x^{\star})+\frac{\gamma}{2}\norm{s_k}^{2},
\qquad\gamma>0\ \text{to be chosen}. \tag{3}
\]

\item \textbf{Bound on the function term.}  
From the descent lemma (Lemma~\ref{lem:descent}) with $y=x_{k+1}$ and $x=x_k$,
\[
f(x_{k+1})\le f(x_k)+\inner{\nabla f(x_k)}{s_{k+1}}+\frac{L}{2}\norm{s_{k+1}}^{2}. \tag{4}
\]

\item \textbf{Express $\inner{\nabla f(x_k)}{s_{k+1}}$ via the update.}  
From (1) we have $-\alpha\nabla f(x_k)=e_{k+1}-e_k-\beta s_k$. Taking inner product with $s_{k+1}=e_{k+1}-e_k$ yields
\[
\begin{aligned}
\inner{\nabla f(x_k)}{s_{k+1}}
&=-\frac{1}{\alpha}\inner{e_{k+1}-e_k-\beta s_k}{e_{k+1}-e_k}\\
&=-\frac{1}{\alpha}\norm{s_{k+1}}^{2}+\frac{\beta}{\alpha}\inner{s_k}{s_{k+1}}. \tag{5}
\end{aligned}
\]

\item \textbf{Insert (5) into (4).}  
\[
\begin{aligned}
f(x_{k+1})&\le f(x_k)-\frac{1}{\alpha}\norm{s_{k+1}}^{2}
+\frac{\beta}{\alpha}\inner{s_k}{s_{k+1}}+\frac{L}{2}\norm{s_{k+1}}^{2}\\
&= f(x_k)-\underbrace{\Bigl(\frac{1}{\alpha}-\frac{L}{2}\Bigr)}_{=:c_1}\norm{s_{k+1}}^{2}
+\frac{\beta}{\alpha}\inner{s_k}{s_{k+1}}. \tag{6}
\end{aligned}
\]

\item \textbf{Add the kinetic term.}  
From (3) we have
\[
\Phi_{k+1}=f(x_{k+1})-f(x^{\star})+\frac{\gamma}{2}\norm{s_{k+1}}^{2}.
\]
Using (6) and cancelling $f(x^{\star})$,
\[
\Phi_{k+1}\le f(x_k)-f(x^{\star})
-\bigl(c_1-\tfrac{\gamma}{2}\bigr)\norm{s_{k+1}}^{2}
+\frac{\beta}{\alpha}\inner{s_k}{s_{k+1}}. \tag{7}
\]

\item \textbf{Upper bound the cross term.}  
Apply Young’s inequality $2\inner{a}{b}\le \eta\norm{a}^{2}+\eta^{-1}\norm{b}^{2}$ with $\eta>0$ to $\inner{s_k}{s_{k+1}}$:
\[
\inner{s_k}{s_{k+1}}\le \frac{\eta}{2}\norm{s_k}^{2}+\frac{1}{2\eta}\norm{s_{k+1}}^{2}. \tag{8}
\]

\item \textbf{Choose $\eta$ to cancel the $s_{k+1}$ term.}  
Plug (8) into (7):
\[
\Phi_{k+1}\le f(x_k)-f(x^{\star})
-\Bigl(c_1-\tfrac{\gamma}{2}-\frac{\beta}{2\alpha\eta}\Bigr)\norm{s_{k+1}}^{2}
+\frac{\beta\eta}{2\alpha}\norm{s_k}^{2}. \tag{9}
\]

Set $\gamma$ and $\eta$ such that
\[
c_1-\frac{\gamma}{2}-\frac{\beta}{2\alpha\eta}=0
\quad\Longleftrightarrow\quad
\gamma=2c_1-\frac{\beta}{\alpha\eta}. \tag{10}
\]

Because $\gamma>0$, we need $2c_1>\frac{\beta}{\alpha\eta}$. Choose
\[
\eta:=\frac{\beta}{\alpha(2c_1-\gamma)}\quad\text{with}\quad
\gamma:=c_1-\frac{L}{2}\beta. \tag{11}
\]
Observe that $c_1=\frac{1}{\alpha}-\frac{L}{2}$, thus $\gamma>0$ precisely when $\alpha<\frac{2}{L+\mu}$ (used later).

\item \textbf{Resulting recursion for $\Phi_k$.}  
With the choice (10) the $s_{k+1}$ term disappears and (9) reduces to
\[
\Phi_{k+1}\le f(x_k)-f(x^{\star})+\frac{\beta\eta}{2\alpha}\norm{s_k}^{2}
= \Phi_k-\bigl(1-\tfrac{\beta\eta}{2\alpha}\bigr)\bigl[f(x_k)-f(x^{\star})\bigr]. \tag{12}
\]

\item \textbf{Link function suboptimality to the Lyapunov function.}  
From strong convexity,
\[
f(x_k)-f(x^{\star})\ge\frac{\mu}{2}\norm{e_k}^{2}. \tag{13}
\]
Moreover, by definition of $\Phi_k$,
\[
\Phi_k\ge f(x_k)-f(x^{\star}). \tag{14}
\]

\item \textbf{Derive a linear contraction.}  
Combine (12)–(14):
\[
\Phi_{k+1}\le\Bigl(1-\frac{\mu}{2}\bigl(1-\tfrac{\beta\eta}{2\alpha}\bigr)\Bigr)\Phi_k
=: \rho\,\Phi_k,
\]
where
\[
\rho:=1-\frac{\mu}{2}\Bigl(1-\frac{\beta\eta}{2\alpha}\Bigr)\in(0,1)
\]
by the parameter restrictions \eqref{ass:params}. \hfill\qedsymbol
\end{enumerate}

\section*{Rate Derivation (With Constants)}
From the contraction $\Phi_{k+1}\le\rho\Phi_k$ we obtain by induction
\[
\Phi_k\le\rho^{k-1}\Phi_1,\qquad k\ge1.
\]
Since $\Phi_k\ge f(x_k)-f(x^{\star})$, the function values satisfy
\[
f(x_k)-f(x^{\star})\le\rho^{k-1}\Phi_1.
\]
Using strong convexity again,
\[
\norm{x_k-x^{\star}}^{2}\le\frac{2}{\mu}\bigl[f(x_k)-f(x^{\star})\bigr]
\le\frac{2\Phi_1}{\mu}\,\rho^{\,k-1},
\]
which yields the stated $O(\rho^{k})$ linear rate.

\section*{Special Cases}
\begin{itemize}
\item \textbf{Exact quadratic $f(x)=\tfrac12 x^{\top}Qx-b^{\top}x$.}  
The error recursion becomes linear: $e_{k+1}= (I-\alpha Q+\beta I)e_k-\beta e_{k-1}$. The spectral radius of the companion matrix is exactly $\rho$ derived above; the bound is tight.
\item \textbf{Gradient descent ($\beta=0$).}  
The Lyapunov function reduces to $f(x_k)-f(x^{\star})$ and $\rho=1-\alpha\mu$, recovering the classical $O((1-\alpha\mu)^k)$ rate.
\item \textbf{Optimal parameters.}  
Setting $\alpha^{\star}=4/(L+\mu)$ and $\beta^{\star}=((\sqrt{\kappa}-1)/(\sqrt{\kappa}+1))^{2}$ yields
\[
\rho^{\star}= \frac{\sqrt{\kappa}-1}{\sqrt{\kappa}+1},
\]
the optimal linear rate for the Heavy–Ball method on quadratics.
\end{itemize}

\end{document}